% !TEX root = ../main.tex
% Figure: RCWA unit cell discretization into layers
% Chapter: ch13 - Rigorous Coupled-Wave Analysis
% Description: Staircase approximation of arbitrary profiles

\begin{figure}[htbp]
  \centering
  \resizebox{\textwidth}{!}{%
  \begin{tikzpicture}[scale=1.15]
    % Original continuous profile (slanted grating)
    \begin{scope}[shift={(-4,0)}]
      % Substrate
      \fill[gray!30] (-2,-1) rectangle (2,0);
      \draw[thick] (-2,0) -- (2,0);
      \node[right] at (2.1,-0.5) {衬底$n_s$};
      
      % Slanted grating teeth
      \foreach \x in {-1.5,-0.75,...,1.5} {
        \fill[blue!20] (\x,0) -- (\x+0.3,0) -- (\x+0.5,1.5) -- (\x+0.2,1.5) -- cycle;
        \draw[thick, blue] (\x,0) -- (\x+0.3,0) -- (\x+0.5,1.5) -- (\x+0.2,1.5) -- cycle;
      }
      
      % Superstrate (air)
      \fill[blue!5] (-2,1.5) rectangle (2,2.5);
      \draw[thick] (-2,1.5) -- (2,1.5);
      \node[right] at (2.1,2) {超strate$n_0$};
      
      % Incident plane wave
      \draw[->, very thick, red] (-2.5,1) -- (-1.8,1.3);
      \node[left, red] at (-2.5,1) {入射光};
      
      % Period indication
      \draw[<->, thick] (-1.5,-0.3) -- node[below] {$\Lambda$} (-0.75,-0.3);
      
      % Height indication
      \draw[<->, thick] (2.3,0) -- node[right] {$h$} (2.3,1.5);
      
      % Tilt angle
      \draw[->] (-1.5,0) arc (90:75:0.8);
      \node[right] at (-1.3,0.6) {\scriptsize $\theta_t$};
      
      \node[above, font=\bfseries] at (0,2.8) {(a) 真实倾斜光栅};
      
      % Fill factor label
      \node[align=center, font=\small] at (0,-0.8) {占空比$f=w/\Lambda$\\锥角$\theta_t$};
    \end{scope}
    
    % Discretized staircase approximation
    \begin{scope}[shift={(2,0)}]
      % Substrate
      \fill[gray!30] (-2,-1) rectangle (2,0);
      \draw[thick] (-2,0) -- (2,0);
      \node[right] at (2.1,-0.5) {衬底$n_s$};
      
      % Staircase layers (10 slices)
      \foreach \y [count=\i] in {0.15,0.3,...,1.5} {
        \pgfmathsetmacro{\width}{1.5 - \i*0.1}
        \fill[green!\i0!blue!20] (-1.5,\y-0.15) rectangle ++(\width,0.15);
        \draw[dashed, gray] (-2,\y) -- (2,\y);
      }
      
      % Layer numbering
      \node[left, font=\small] at (-2.1,0.75) {$\ell=1$};
      \node[left, font=\small] at (-2.1,1.5) {$\ell=N_L$};
      
      % Effective index in each layer
      \node[right, font=\footnotesize] at (1.2,0.5) {$n_{\mathrm{eff}}^{(1)}$};
      \node[right, font=\footnotesize] at (1.2,1.0) {$n_{\mathrm{eff}}^{(2)}$};
      \node[right, font=\footnotesize] at (1.2,1.4) {$n_{\mathrm{eff}}^{(N_L)}$};
      
      % Superstrate
      \fill[blue!5] (-2,1.5) rectangle (2,2.5);
      \draw[thick] (-2,1.5) -- (2,1.5);
      \node[right] at (2.1,2) {超strate$n_0$};
      
      \node[above, font=\bfseries] at (0,2.8) {(b) 阶梯近似分层};
      
      % Number of layers
      \draw[<->, thick] (2.3,0) -- node[right] {$N_L=10$} (2.3,1.5);
    \end{scope}
    
    % Connecting arrow showing approximation concept
    \draw[->, ultra thick, orange, bend left=15] (-1.5,2.5) to[out=20,in=160] 
      node[above, align=center, font=\small] {Li 因子化 + 傅里叶谐波展开\\\\tiny 每层内介电常数$\varepsilon(x,y)$分段恒定} (0,2.5);
    
    % Bottom: Fourier series representation
    \node[draw, align=center, font=\small] at (0,-1.5) {
      第$\ell$层介电函数的傅里叶展开:\\
      $\varepsilon^{(\ell)}(x) = \sum_{m=-\infty}^{\infty}\hat{\varepsilon}_m^{(\ell)}e^{i m G x}$,\\
      其中$G=2\pi/\Lambda$,截断至$m=\pm N_h$个谐波\\[4pt]
      S-matrix 算法保证数值稳定性
    };
    
    % Convergence note
    \node[align=left, font=\footnotesize, anchor=south west] at (-3.5,-2.5) {
      \textbf{收敛准则}:\\
      $\bullet$ 谐波数$N_h\geq 21$保证$<1\%$误差\\
      $\bullet$ 分层数$N_L$:陡峭侧壁需更细划分\\
      $\bullet$ Li 规则处理不连续面避免 Gibbs 震荡
    };
  \end{tikzpicture}%
  }
  \caption{RCWA 中常用的阶梯离散化策略。(a) 上图示意具有倾斜侧壁的真实周期结构；(b) 下图示意将其沿纵向切分为 $N_L$ 个矩形薄层，并在每层内用有限谐波数 $N_h$ 展开介电函数。配合界面连续条件和散射矩阵递推，可得到整体的反射、透射与模式耦合信息。}
  \label{fig:ch13_rcwa_discretization}
\end{figure}
